\section{支撑材料文件列表}
\begin{longtable}{p{0.37\textwidth}p{0.56\textwidth}}
\caption{支撑材料文件及用途}\label{tab:support-files}\\
\toprule
文件 & 用途 \\
\midrule
\endfirsthead
\toprule
文件 & 用途 \\
\midrule
\endhead
\path{main.tex}、\path{shumo-cumcm2025.cls} & 论文的 LaTeX 唯一正文源和 2025 电子版格式类文件 \\
\path{solve_case001.py} & 复现问题一至问题三，并审计问题四冻结路线的条数、总长度和最小曲率半径 \\
\path{make_figures.py} & 从冻结数值和路线数据生成正文所用矢量 PDF 图 \\
\path{data/problem4_final_metrics.json} & 问题四最终指标、敏感性结果和结论边界 \\
\path{data/problem4_route_station_matrix.csv} & $61$ 条路线在 $11$ 个站位上的横向坐标、单线长度和最小曲率半径 \\
\path{figures/*.pdf} & 正文引用的矢量图形 \\
\bottomrule
\end{longtable}

\section{问题四路线摘要}
表~\ref{tab:route-summary} 给出每条测线在南端、中部和北端的横向坐标，以及显式几何长度和最小曲率半径。完整 $11$ 站位数据见支撑材料 CSV 文件。

\begin{longtable}{C{1.4cm}C{1.65cm}C{1.65cm}C{1.65cm}C{2.15cm}C{2.15cm}}
\caption{问题四 $61$ 条推荐测线摘要}\label{tab:route-summary}\\
\toprule
测线 & $x(0)$ /海里 & $x(2.5)$ /海里 & $x(5)$ /海里 & 长度/m & 最小曲率半径/m \\
\midrule
\endfirsthead
\multicolumn{6}{c}{表~\thetable\ （续）}\\
\toprule
测线 & $x(0)$ /海里 & $x(2.5)$ /海里 & $x(5)$ /海里 & 长度/m & 最小曲率半径/m \\
\midrule
\endhead
\midrule
\multicolumn{6}{r}{续下页}\\
\endfoot
\bottomrule
\endlastfoot
C-001 & 0.0222 & 0.0318 & 0.0792 & 9260.95 & 29655.46 \\
C-002 & 0.0591 & 0.0786 & 0.1503 & 9262.19 & 14458.83 \\
C-003 & 0.0952 & 0.1252 & 0.2215 & 9263.99 & 9527.04 \\
C-004 & 0.1307 & 0.1714 & 0.2928 & 9266.40 & 7119.43 \\
C-005 & 0.1656 & 0.2174 & 0.3641 & 9269.42 & 5706.11 \\
C-006 & 0.2000 & 0.2631 & 0.4354 & 9273.08 & 4809.65 \\
C-007 & 0.2339 & 0.3087 & 0.5068 & 9277.38 & 4126.31 \\
C-008 & 0.2675 & 0.3540 & 0.5781 & 9282.34 & 3611.11 \\
C-009 & 0.3007 & 0.3993 & 0.6495 & 9287.96 & 3220.54 \\
C-010 & 0.3336 & 0.4445 & 0.7208 & 9294.24 & 2915.70 \\
C-011 & 0.3662 & 0.4896 & 0.7921 & 9301.18 & 2662.63 \\
C-012 & 0.3987 & 0.5347 & 0.8633 & 9308.74 & 2447.99 \\
C-013 & 0.4310 & 0.5799 & 0.9344 & 9316.95 & 2271.09 \\
C-014 & 0.4632 & 0.6251 & 1.0054 & 9325.74 & 2118.58 \\
C-015 & 0.4953 & 0.6704 & 1.0765 & 9335.16 & 1987.89 \\
C-016 & 0.5274 & 0.7158 & 1.1473 & 9345.10 & 1873.22 \\
C-017 & 0.5596 & 0.7614 & 1.2181 & 9355.61 & 1770.72 \\
C-018 & 0.5918 & 0.8072 & 1.2886 & 9366.59 & 1680.94 \\
C-019 & 0.6242 & 0.8533 & 1.3592 & 9378.07 & 1602.69 \\
C-020 & 0.6568 & 0.8997 & 1.4294 & 9389.90 & 1533.92 \\
C-021 & 0.6895 & 0.9464 & 1.4996 & 9402.18 & 1472.09 \\
C-022 & 0.7226 & 0.9934 & 1.5694 & 9414.71 & 1408.26 \\
C-023 & 0.7560 & 1.0408 & 1.6393 & 9427.59 & 1349.98 \\
C-024 & 0.7897 & 1.0888 & 1.7087 & 9440.59 & 1305.61 \\
C-025 & 0.8239 & 1.1372 & 1.7781 & 9453.84 & 1265.22 \\
C-026 & 0.8586 & 1.1861 & 1.8470 & 9467.09 & 1221.12 \\
C-027 & 0.8939 & 1.2356 & 1.9160 & 9480.48 & 1180.72 \\
C-028 & 0.9298 & 1.2858 & 1.9844 & 9493.72 & 1143.90 \\
C-029 & 0.9663 & 1.3367 & 2.0528 & 9506.95 & 1115.05 \\
C-030 & 1.0037 & 1.3884 & 2.1207 & 9519.88 & 1085.75 \\
C-031 & 1.0419 & 1.4408 & 2.1885 & 9532.72 & 1058.61 \\
C-032 & 1.0810 & 1.4941 & 2.2558 & 9544.98 & 1029.64 \\
C-033 & 1.1211 & 1.5483 & 2.3230 & 9556.99 & 1002.09 \\
C-034 & 1.1623 & 1.6035 & 2.3896 & 9568.35 & 971.74 \\
C-035 & 1.2047 & 1.6600 & 2.4561 & 9579.22 & 920.25 \\
C-036 & 1.2485 & 1.7175 & 2.5220 & 9589.19 & 871.86 \\
C-037 & 1.2938 & 1.7763 & 2.5879 & 9598.51 & 826.52 \\
C-038 & 1.3408 & 1.8363 & 2.6531 & 9606.76 & 777.45 \\
C-039 & 1.3896 & 1.8977 & 2.7182 & 9614.18 & 730.90 \\
C-040 & 1.4403 & 1.9608 & 2.7826 & 9620.17 & 690.44 \\
C-041 & 1.4933 & 2.0257 & 2.8471 & 9625.01 & 658.97 \\
C-042 & 1.5486 & 2.0922 & 2.9106 & 9628.35 & 618.34 \\
C-043 & 1.6065 & 2.1605 & 2.9742 & 9630.38 & 579.81 \\
C-044 & 1.6674 & 2.2307 & 3.0370 & 9630.46 & 548.43 \\
C-045 & 1.7314 & 2.3031 & 3.0997 & 9628.85 & 519.78 \\
C-046 & 1.7991 & 2.3783 & 3.1616 & 9625.25 & 486.09 \\
C-047 & 1.8707 & 2.4558 & 3.2235 & 9619.62 & 458.59 \\
C-048 & 1.9468 & 2.5359 & 3.2844 & 9611.52 & 432.44 \\
C-049 & 2.0279 & 2.6188 & 3.3454 & 9601.19 & 407.59 \\
C-050 & 2.1153 & 2.7046 & 3.4055 & 9588.30 & 380.71 \\
C-051 & 2.2093 & 2.7945 & 3.4655 & 9572.75 & 357.72 \\
C-052 & 2.3106 & 2.8881 & 3.5246 & 9554.21 & 336.46 \\
C-053 & 2.4206 & 2.9856 & 3.5837 & 9532.90 & 319.56 \\
C-054 & 2.5405 & 3.0872 & 3.6419 & 9509.49 & 298.12 \\
C-055 & 2.6719 & 3.1932 & 3.7000 & 9483.98 & 278.33 \\
C-056 & 2.8167 & 3.3048 & 3.7571 & 9456.18 & 263.57 \\
C-057 & 2.9774 & 3.4226 & 3.8142 & 9427.20 & 251.71 \\
C-058 & 3.1569 & 3.5464 & 3.8704 & 9400.22 & 236.82 \\
C-059 & 3.3591 & 3.6768 & 3.9265 & 9376.87 & 225.98 \\
C-060 & 3.5892 & 3.8143 & 3.9633 & 9346.04 & 261.87 \\
C-061 & 3.8534 & 3.9594 & 4.0000 & 9298.75 & 503.01 \\

\end{longtable}

\section{核心复现程序}
以下程序使用标准 Python 3 库即可运行，复现问题一至问题三的数值结果，并检查问题四最终路线包的关键一致性。
\lstinputlisting[caption={问题一至问题四核心结果复现与审计程序}]{solve_case001.py}

\section{绘图程序}
\lstinputlisting[caption={正文矢量图生成程序}]{make_figures.py}
